\begin{question}

A beam of silver atoms is created by heating a vapor in an oven to
1000C, and selecting atoms with a velocity close to the mean of the
thermal distribution. The beam moves through a one-meter long magnetic
field with a vertical gradient 10 Bohr magnetons, and impinges a
screen one meter downstream of the end of the magnet. Assuming the
silver atom has spin $\frac 1 2$ with a magnetic moment of one Bohr
magneton, find the separation distance in millimeters of the two
states on the screen.

\end{question}


\begin{given}

\begin{equation}
\frac{\partial}{\partial x}B = 10 \frac T m
\end{equation}

\begin{equation}
T_0 = 1000C
\end{equation}

\begin{equation}
  x_0 = 1m
\end{equation}

\begin{equation}
  \mu = 1 \frac J T
\end{equation}

\end{given}

From basic Newtonian physics, we have

\begin{equation}
\begin{split}
  S &= \frac 1 2 a t^2 \\
  &= \frac 1 2 a \frac{x_0^2}{v^2} \\
  &= \frac 1 2 \frac F m \frac{x_0^2}{v^2} \\
\end{split}
\end{equation}

This leaves us with an equation that requires a Force and a velocity.
The temperature provides the energy which can be employed to solve for
the velocity

\begin{equation}
\begin{split}
  E &= \frac 3 2 K_bT \\
  &= \frac{p^2}{2m} \\
  &= \frac 1 2 m v^2
\end{split}
\end{equation}

Solving for v, we find that

\begin{equation}
  v = \left(\frac{3 K_b T}{m}\right)^\frac{1}{2}
\end{equation}

Substituting our expression for the velocity in, we find that:

\begin{equation}
\begin{split}
  S &= \frac 1 2 \frac F m \frac{x_0^2}{v^2} \\
  &= \frac 1 2 \frac F m x_0^2 \frac{m}{3K_bT} \\
  &= \frac{F x_0^2}{6K_bT}
\end{split}
\end{equation}

The force ($F_0$) can be determined by using the magnetic field
gradient:

\begin{equation}
  F = \nabla(\vec{\mu} \cdot \vec{B}) = 10 \cdot \mu_B \approx 9.275\cdot 10^{-23} N
\end{equation}

Armed with the force and velocity we can now solve for
the displacement:

\begin{equation}
  S = \frac{F_0 x_0^2}{6K_bT} \approx 0.88mm
\end{equation}
